\section{Integer bands from traceable subdivisions of $K_4$}
\label{sec:rhost-k4-bands}

A graph is traceable if it has a Hamilton path.
Every subdivision below retains one edge of $K_4$ without subdivision.
This is necessary for the stated spanning-path construction.
A subdivision with an internal vertex on all six arms has no Hamilton path.
Indeed, such a path would delete three edges while reducing all four branch degrees.
Each deleted edge would then meet at most one branch vertex.

\begin{lemma}\label{lem:rhost-k4-bands}
For each integer $N\ge1$, there are $2(N+1)$ traceable subdivisions of $K_4$ with these properties:
\begin{enumerate}
\item Every cycle has length at least $2N+6$.
\item Every bicyclic core has at least $7N+13$ edges.
\item Each whole subdivision has at least $14N+15$ edges and at most $15N+14$ vertices.
\item Their cycle lengths include every integer in $[2N+6,14N+10]$ except possibly $4N+8$.
\end{enumerate}
Consequently, adding $2N+4$ individual cycles covers all lengths $3,\ldots,14N+10$.
The total cycle rank of these blocks is $8N+10$.
Their total number of vertices is at most $(61N^2+141N+94)/2$.
\end{lemma}

\begin{proof}
Label the vertices of $K_4$ by $a,b,c,d$.
List its six arms in the order
\[
 (ab,ac,ad,bc,bd,cd).
\]
For each integer $t\in\{0,1,\ldots,N\}$, define two subdivisions by the following arm lengths:
\begin{align*}
 A_t&=(8N-2t+3,\ 1,\ N-t+3,\ 2N+2t+3,\ N+t+3,\ 3N-t+3),\\
 B_t&=(1,\ N-t+2,\ N-t+3,\ 3N-t+3,\ 3N+3t+3,\ 6N+3).
\end{align*}
All arm lengths are positive integers.
The arm $ac$ of $A_t$ and the arm $ab$ of $B_t$ each have length one.
Every other arm has length at least two.

The six selected cycle lengths in the following table each run through an arithmetic progression with step two.
The interval gives the two endpoint values, irrespective of direction.
\[
\begin{array}{c|cc|cc}
 \text{band}&\text{cycle in }A_t&\text{length interval}&
              \text{cycle in }B_t&\text{length interval}\\\hline
 0&acd&[2N+7,4N+7]&abc&[2N+6,4N+6]\\
 1&acbd&[4N+10,6N+10]&abd&[4N+7,6N+7]\\
 2&bcd&[6N+9,8N+9]&acd&[6N+8,8N+8]\\
 3&abd&[8N+9,10N+9]&abcd&[8N+10,10N+10]\\
 4&abdc&[10N+10,12N+10]&abdc&[10N+9,12N+9]\\
 5&abcd&[12N+12,14N+12]&bcd&[12N+9,14N+9]
\end{array}
\]
The unlisted cycles have lengths $10N+7$ in $A_t$ and $8N+11$ in $B_t$.
These formulas show that every cycle has length at least $2N+6$.
Each band supplies the two opposite parities.
The endpoint intervals overlap across all interfaces except the single possible gap $4N+8$.
Their union therefore contains $[2N+6,14N+10]\setminus\{4N+8\}$.

The total arm lengths are
\[
 E(A_t)=15N-t+16,\qquad E(B_t)=14N+15.
\]
A bicyclic core in a subdivision of $K_4$ consists of five complete arms.
For $A_t$, the longest arm is $ab$, so every such core has at least $7N+t+13$ edges.
For $B_t$, the longest arm is $cd$, tied by $bd$ when $t=N$.
Every such core therefore has at least $8N+12\ge7N+13$ edges.
The whole subdivisions have cycle rank three and $E-2$ vertices.
This proves the edge and vertex bounds.

It remains to verify traceability.
In $A_t$, omit the edge $ac$.
Delete the edge at $b$ on the arm $ab$, and the edge at $d$ on the arm $cd$.
The remaining graph is the spanning path whose branch order is
\[
 a,d,b,c.
\]
Its endpoints are the internal vertices adjacent to $b$ on $ab$ and to $d$ on $cd$.
In $B_t$, omit $ab$.
Delete the edge at $c$ on $ac$, and the edge at $d$ on $bd$.
The remaining graph is a spanning path with branch order
\[
 a,d,c,b.
\]
Its endpoints are the internal vertices adjacent to $c$ on $ac$ and to $d$ on $bd$.
Every arm used to specify these endpoints has an internal vertex.

Add one individual cycle of each length $3,\ldots,2N+5$, and one of length $4N+8$.
This uses $2N+4$ cycles and fills the prefix and the possible gap.
The cycle-rank sum is
\[
 3\cdot2(N+1)+(2N+4)=8N+10.
\]
The $K_4$ subdivisions together have
\[
 \sum_{t=0}^{N}\bigl(15N-t+14+14N+13\bigr)
 =\frac{57N^2+111N+54}{2}
\]
vertices.
The individual cycles add $2N^2+15N+20$ vertices.
Adding these expressions gives the claimed total.
\end{proof}

For a target interval $3,\ldots,T$, set $N=\lceil T/14\rceil$.
The lemma gives total block rank at most $4T/7+18$.
All blocks together use $O(T^2)$ vertices.
If $d=7N$, the whole-edge threshold is at least $2d$, every bicyclic core has at least $d$ edges,
and every cycle has at least $2d/7$ edges.
These inequalities are the local density inputs for packing the blocks into a random host.
The lemma itself does not assert a random-host embedding theorem.

The exact formulas and the Hamilton paths are checked by
\nolinkurl{check_k4_random_host_blocks.py}.
